% Transcription — fonds Alexandre Grothendieck, shelfmark 35, batch 1 (pages 1-5).
% Established from the facsimile archives/batches/35/batch-01.pdf (a mirror of
% https://grothendieck.umontpellier.fr/35.pdf), per the skill
% transcribe-grothendieck.
% The plan of SGA 7. Page 1 is a typed table of exposés I to IX which the author
% has gone over in ink, renumbering, retitling, and writing a speaker's name
% against each line in the left margin. Page 2 continues X to XXII in his hand,
% with D or K against each — initials of the two who were to write them, which
% the page does not develop and which are not developed here. Page 3 is the tail
% of a bibliography, [17] to [23], struck through. Page 4 is a leaf of typescript
% on bitorseurs, paginated (11bis) in its own numbering. Page 5 carries three
% dependency graphs of the exposés: the first is crossed out by four large
% strokes and is not redrawn, since its arrows run under the cancellation and are
% no longer all legible; the two below it, its replacement, are given.
% Pass: Opus 5 (claude-opus-5), 2026-08-16 — first pass, unchecked against the pages by a human.
\documentclass[11pt,a4paper]{article}
% Le préambule commun à toutes les transcriptions du fonds Grothendieck.
%
% Il tient en une idée : **l'incertitude se compose, elle ne se résout pas.**
% Une transcription qui lisse un mot douteux a détruit la seule chose pour
% laquelle on la faisait — un lecteur doit pouvoir savoir, à chaque endroit,
% ce qui était sur le feuillet et ce qui a été ajouté. Les six macros qui
% suivent sont donc l'appareil critique, et non de la décoration.
%
% Les mêmes commandes sont reconnues par `scripts/render.mjs`, qui produit la
% vue de lecture du site. Une macro ajoutée ici doit l'être aussi là-bas,
% faute de quoi le rendu échouera bruyamment — ce qui est voulu : un rendu
% silencieusement faux est pire qu'un rendu absent.

\ProvidesPackage{grothendieck}[2026/08/08 Transcriptions du fonds Grothendieck]

\usepackage{amsmath,amssymb,amsthm}
\usepackage{mathrsfs}
\usepackage{stmaryrd}
% Les diagrammes commutatifs sont partout dans ce fonds : c'est la langue
% même de la géométrie algébrique de Grothendieck, et une transcription qui
% les rendrait en prose ne transcrirait rien.
\usepackage{tikz-cd}
% Français par défaut — c'est la langue des feuillets — mais l'anglais est
% chargé aussi : la traduction et le résumé sont des documents anglais, et la
% typographie française (espaces avant les deux-points) y serait une faute.
% Ces documents basculent par \selectlanguage{english} en tête de corps.
\usepackage[english,french]{babel}
\usepackage{fontspec}
\usepackage{geometry}
\geometry{a4paper,margin=25mm,marginparwidth=28mm}
\usepackage{marginnote}
\usepackage{xcolor}
% ulem, not soul. `soul`'s \st and \ul are built on a character-by-character
% scan that XeLaTeX's font machinery breaks: the document fails with an
% undefined control sequence rather than degrading. ulem does the same two jobs
% and is Unicode-engine safe. `normalem` stops it hijacking \emph, which the
% transcriptions use for Grothendieck's own emphasis.
\usepackage[normalem]{ulem}
\usepackage{hyperref}
\hypersetup{colorlinks=true,linkcolor=black,urlcolor=black,pdfborder={0 0 0}}

% --- Métadonnées du lot -----------------------------------------------------
% Reprises telles quelles par le script de rendu, qui les affiche en tête de
% la vue de lecture. Elles fixent ce que le fichier prétend couvrir : sans
% elles, une transcription flotte sans point d'attache dans un fonds de seize
% mille feuillets.

\newcommand{\folder}[1]{\gdef\grothFolder{#1}}
\newcommand{\batch}[1]{\gdef\grothBatch{#1}}
\newcommand{\leaves}[2]{\gdef\grothFirst{#1}\gdef\grothLast{#2}}
\newcommand{\foldertitle}[1]{\gdef\grothFoldertitle{#1}}
\gdef\grothFolder{?}\gdef\grothBatch{?}\gdef\grothFirst{?}\gdef\grothLast{?}\gdef\grothFoldertitle{}

% --- Le résumé --------------------------------------------------------------
%
% Il ouvre la lecture modernisée, sous le titre « Résumé », et il est composé
% en petit. Ce n'est pas un ornement : c'est le seul passage du document qui
% ne soit pas des mathématiques, et un lecteur doit pouvoir le reconnaître
% avant de l'avoir lu — puis le sauter s'il connaît déjà le sujet, ou s'y
% arrêter s'il ne le connaît pas. Le filet qui le suit marque la reprise du
% corps.

\newenvironment{resume}
  {\small\setlength{\parskip}{0.45em}}
  {\par\medskip\hrule\medskip}

% --- Mots-clés ---------------------------------------------------------------
%
% En anglais, en clôture du résumé de la lecture modernisée : le vocabulaire
% moderne sous lequel chercher ce que les feuillets construisent. Le manifeste
% du site (`npm run manifest`) les extrait pour étiqueter la cote entière —
% cette ligne est la source unique des tags, il n'y a pas de fichier à côté.
\newcommand{\keywords}[1]{\par\smallskip\noindent
  {\footnotesize\textsc{Keywords} --- \textit{#1}}\par}

% --- L'appareil critique ----------------------------------------------------

% Le numéro de feuillet, en marge. C'est l'ancre qui permet de régler un
% désaccord sur une formule en regardant *un* feuillet plutôt qu'une chemise
% de six cents. Le site s'en sert aussi pour tourner les pages du fac-similé
% quand on fait défiler la transcription.
\newcommand{\leaf}[1]{%
  \marginnote{\footnotesize\color{gray!70}\textsf{#1}}%
  \label{leaf:#1}\ignorespaces}

% Illisible : on ne devine pas. Un mot inventé qui se lit comme les autres est
% le pire résultat possible.
\newcommand{\ill}{\textcolor{red!60!black}{[\,\ldots\,]}}

% Lecture douteuse : le mot est donné, et signalé comme incertain.
\newcommand{\uncertain}[1]{\uline{#1}}

% Ajout de l'éditeur — une lettre manquante, un mot sous-entendu.
\newcommand{\add}[1]{\textcolor{blue!50!black}{[#1]}}

% Biffé par Grothendieck lui-même. Ce qu'il a rayé fait partie du document :
% une rature dit souvent où la pensée a changé de direction.
\newcommand{\struck}[1]{\sout{#1}}

% Note du transcripteur, sur l'état matériel du feuillet.
\newcommand{\note}[1]{\par\smallskip{\small\color{gray!60!black}\textit{[#1]}}\par\smallskip}

% Note marginale de Grothendieck, distincte du corps du texte.
\newcommand{\marginal}[1]{\par\smallskip\noindent\hspace{1em}%
  {\small\color{gray!60!black}#1}\par\smallskip}

% --- Titre ------------------------------------------------------------------

\AtBeginDocument{%
  \begin{center}
    {\large\bfseries Fonds Alexandre Grothendieck --- cote n\textsuperscript{o}~\grothFolder}\\[2pt]
    {\itshape\grothFoldertitle}\\[4pt]
    {\small feuillets \grothFirst--\grothLast\ (lot \grothBatch)}\\[2pt]
    {\footnotesize\color{gray!60!black}Transcription établie d'après le fac-similé numérisé
     par l'Université de Montpellier.\\
     Les lectures incertaines sont soulignées, les illisibles notées [\,\ldots\,],
     les ajouts entre crochets.}
  \end{center}
  \vspace{1em}}

\endinput


\folder{35}
\batch{1}
\pages{1}{5}
\dating{[vers 1967-1973]}
\watermark{Édition de démonstration}
\foldertitle{Plan [SGA 7] : notes manuscrites (s.d.), tapuscrit annoté (s.d.)}

\begin{document}

\page{1}

En haut à droite, à la main : \uncertain{dernier en G}.

\section*{SGA 7 --- Groupe de Monodromie en Géométrie Algébrique}

\note{le tapuscrit donne les neuf premiers exposés. La colonne ajoutée à la main
dans la marge de gauche porte, en regard de chaque ligne, celui qui devait la
faire ; on la reporte ici en tête de chaque entrée, après le numéro. « Gr. » est
l'auteur lui-même. Les corrections dans le corps sont de la même encre que
cette colonne}

\begin{itemize}
\item[I --- Gr.] Généralités sur le groupe de monodromie locale. Le théorème de
monodromie par voie arithmétique.

\item[II --- Gr.] Cas d'actions résolubles d'un groupe fondamental (variante du
\struck{groupe} théorème de monodromie).

\item[III --- Gr.] Le théorème de monodromie par voie géométrique.

\item[IV --- Gr.] Propriétés générales des groupes de cycles évanescents. Cas
des points singuliers isolés de la fibre spéciale.

\item[V --- Gr.] Groupe de monodromie opérant sur le groupe fondamental :
théorèmes d'action modérée et d'action essentiellement modérée.

\item[\struck{VI} Arr.] Application : présentation finie des groupes
fondamentaux géométriques

\item[VI] \struck{Modules} Modules formels (\emph{ist bei Deligne})

\item[VII --- Gr.] Biextensions de faisceaux de groupes.

\item[VIII --- Gr.] \add{Extensions et biextensions de schémas en groupes}
\ill\ $\mathbb{G}_{m}$ d'orthogonalité, théorème

\item[IX --- Gr.] Applications au modèle de Néron : \struck{\ill} théorème de
réduction semi-stable\ldots

\item[\struck{X}] \struck{Application : la réduction semi-stable des courbes
algébriques, et la connexité des variétés de modules pour les courbes de genre
$g$.}
\end{itemize}

\marginal{en regard du III : indépendance de l'action \ill\ sur cycles
évanescents}

\marginal{en regard du V, cerclé et suivi d'un point d'interrogation : Mme
Raynaud ?}

\marginal{en regard de l'exposé renuméroté « Arr. » : \uncertain{Marc} Raynaud,
écrit au-dessus de \struck{Mme Raynaud}}

\marginal{en regard du VI « Modules formels » : \struck{Rim}}

\marginal{en regard du dernier, celui qui est entièrement barré : Deligne}

\note{les trois mots entre parenthèses au VI sont en allemand sur la page, et
non en français}

\page{2}

\section*{SGA 7 (suite)}

\begin{itemize}
\item[D --- X] Intersections \ldots

\item[D --- XI] Int.\ complètes

\item[D --- XII] Quadriques

\item[D --- XIII] Points \uncertain{quadratiques} ordinaires. Th.-Lef.\ faible

\item[D --- XIV] Formule Milnor

\item[D --- XV] Picard-Lefschetz

\item[D --- XVI] Picard-Lefschetz

\item[K --- XVII] Pinceaux de Lefschetz : \underline{existence} $/\!/$
\underline{étude cohomologique}

\item[D --- \struck{XVIII} XIX] Th.\ de Nombres, avec \ill\ Hodge--Deligne et
\ill\ Katz

\item[\struck{XX}] Th.\ de Griffiths

\item[K --- XXI] \struck{\ill} Niveau de la coh.\ des intersections complètes

\item[K --- XXII] \struck{\ill} Calcul de la \ill\ mod $p$.
\end{itemize}

\note{une accolade réunit XV et XVI, qui portent le même titre ; une seconde,
dans la marge de gauche, réunit les trois derniers en regard de K}

\emph{Sous un trait :} Exposés transcendants (Lê Dũng Tráng ?)

\emph{Puis, souligné :} $\mathbb{F}_{q}$ \quad $\#$ \quad card

\page{3}

\note{fin d'une bibliographie, dont ne subsistent que les numéros 17 à 23 ; deux
grands traits de plume la barrent en diagonale. Le reste de la page est vierge}

\begin{itemize}
\item[{[17]}] Griffiths, P. --- Yale talk
\item[{[18]}] Grothendieck, A. --- \emph{Inventiones}
\item[{[19]}] \add{Deligne} \struck{Grothendieck}, A. --- S.G.A.\ on monodromy
--- SGA 7
\item[{[20]}] Landman, A. --- Berkeley thesis (1966)
\item[{[21]}] Lieberman, D. --- \emph{Higher Picard varieties}
\item[{[22]}] Schmid, W. --- \emph{On a Conjecture of Langlands}
\item[{[23]}] Wu, H. --- Acta paper
\end{itemize}

\note{au n\textsuperscript{o} 19 le nom de Deligne est écrit au-dessus de celui
de Grothendieck, qui est barré ; le millésime du n\textsuperscript{o} 20 est
repassé sur un chiffre antérieur}

\page{4}

\note{feuillet de tapuscrit, paginé (11bis) dans sa propre numérotation, repris
à la plume ; deux grands traits le barrent en diagonale comme la bibliographie.
Le passage traite des bitorseurs et de leurs rigidifications, et c'est la
matière de l'exposé VII annoncé p.~1}

1.3.3. La terminologie de 1.3.1 \struck{\ill} s'étend aussitôt au cas où l'on a
un bitorseur $E$ sous $(G_{p}, G_{p})$, \struck{\ill} étant entendu qu'une
rigidification du bitorseur $E$ doit être un isomorphisme de $e_{p}(E)$ avec le
bitorseur trivial $G_{d}$, \add{i.e.} \struck{\ill} est défini par une section
$e_{E}$ de $e_{p}(E)$ qui est \underline{centrale}, i.e.\ qui satisfait à la
condition $g \cdot e_{E} = e_{E} \cdot g$ $(g \in G(S))$. On étend de même la
terminologie des birigidifications au cas des bitorseurs.

1.3.4. Revenons au cas où on se donne deux Groupes $P$ et $Q$, et choisissons la
section unité $e_{p}$ pour définir des rigidifications de bitorseurs sur $P$, et
des birigidifications de bitorseurs sur $P \times P$. On a des foncteurs
naturels

\begin{tikzcd}[column sep=large, nodes={font=\scriptsize}]
\underline{\mathrm{Ext}}(P; G) \arrow[r, "i"] & \underline{\mathrm{Bitorsrig}}(P, G) \arrow[r, "\delta"] & \underline{\mathrm{Bitorsbirig}}(P, P; G)
\end{tikzcd}

où les trois expressions écrites désignent respectivement la catégorie des
extensions de \add{Groupes de} $P$ par $G$, la catégorie des bitorseurs sous
$(G_{p}, G_{p})$ rigidifiés pour $e_{p}$, et la catégorie des bitorseurs sous
$(G_{P \times P}, G_{P \times P})$ birigidifiés pour $(e_{p}, e_{p})$. On
définit le foncteur $i$ en associant à l'extension $E$ de $P$ par $G$ le
\struck{\ill} bitorseur correspondant sur $P$ (1.1.6), rigidifié par la section
unité de $E$. On définit le foncteur $\delta$ par la formule

\note{la formule annoncée manque : le feuillet s'arrête là}

\page{5}

\note{le premier graphe, en haut de la page, est barré de quatre grands traits
croisés et remplacé par les deux qui suivent. On ne le redessine pas : ses
flèches passent sous les ratures et leurs pointes ne sont plus toutes assurées,
et un graphe dont une flèche manque affirme une dépendance que personne n'a
écrite}

\subsection*{Théorèmes qualitatifs sur monodromie (et auxiliaires techniques)}

\begin{tikzcd}[column sep=small, row sep=small, nodes={font=\scriptsize}]
& & & \mathrm{I} \arrow[dlll] \arrow[ddlll] \arrow[ddll] \arrow[dddll]\\
\mathrm{II} & & \mathrm{VI} \arrow[d] \arrow[ddl] & \\
\mathrm{III} \arrow[dd] & \mathrm{IV} & \mathrm{VII} \arrow[d] \arrow[dl, dashed] & \\
& \mathrm{V} & \mathrm{VIII} \arrow[dll] & \\
\mathrm{IX} \arrow[ur] & & &
\end{tikzcd}

\note{les quatre traits qui partent de I sont perlés --- une file de points le
long du fût --- et non pleins ; ils sont rendus ici par des flèches ordinaires,
la distinction n'étant pas dans la syntaxe. La flèche entre IX et V porte sa
pointe en haut, du côté de V : elle est donc donnée dans ce sens, qui est
l'inverse des autres arrivées sur IX, et cette asymétrie est celle de la page.
Une flèche en pointillé arrive encore sur IX par la droite, hors du bloc
transcrit, et porte au-dessous un mot barré et illisible}

\emph{Sous le graphe, une accolade et la légende :} exposés VI, VII, VIII.

\subsection*{Théorèmes quantitatifs}

\begin{tikzcd}[column sep=small, row sep=small, nodes={font=\scriptsize}]
\mathrm{XI} \arrow[d] & & \\
\mathrm{XII} \arrow[d] & & \\
\mathrm{XIII} \arrow[d] \arrow[r] & \mathrm{XIV} \arrow[r, dashed] & \mathrm{XVII} \arrow[d]\\
\mathrm{XV} \arrow[d] & & \mathrm{XVIII} \arrow[ddll]\\
\mathrm{XVI} \arrow[d] & & \\
\mathrm{XIX} & &
\end{tikzcd}

\emph{À part, en bas à droite :} XIX, XX $\leftarrow$ XXI $\leftarrow$ XXII, les
deux premiers cerclés.

\note{le X qui ouvrirait la chaîne est cerclé puis raturé si lourdement qu'il
n'en reste que le contour ; deux traits venus du graphe supérieur y aboutissent
et un pointillé l'atteint par la gauche, mais rien de tout cela n'est assuré et
il est laissé hors du diagramme. Entre XIV et XVIII court un trait terminé par
une pointe et parcouru sur toute sa longueur d'un zigzag serré, qui est la
manière dont l'auteur annule une flèche ailleurs dans le fonds ; il n'est pas
transcrit. Le XII posé plus à gauche, hors chaîne, ne porte aucune flèche}

\marginal{cerclé, à droite du second graphe : \uncertain{quelque chose
d'intérieur !}}

\marginal{NB VI, VII, VIII, XXII \ldots\ tout ce qui est d'un usage de
monodromie, et \ill}

\end{document}
